\scanpage{33}

\deleted{\unsure{Les inclusions}}

on a \uline{loc.} des hom
\[
{}_nA\longrightarrow A,\qquad{}_nA'\longrightarrow A'.
\]
d'où \uline{loc.} \unsure{un hom.}
\diagram{\begin{tikzcd}[ampersand replacement=\&, column sep=large,row sep=large]
\operatorname{Ext}^1(A\mathbin{\overset{L}{\otimes}}A\unsure{'},G_m) \arrow[r] \arrow[dr] \& \operatorname{Ext}^1({}_nA\mathbin{\overset{L}{\otimes}}{}_nA',G_m)\unsure{\to0} \arrow[d]\\
 \& \operatorname{Hom}({}_nA\otimes{}_nA',G_m)
\end{tikzcd}}

\unsure{qui se recollent}, donc \unsure{un hom défini global\textsuperscript{t}}. Donc l'élément $\xi$ envisagé plus haut nous donne
\[
\phi_n^\xi:{}_nA\otimes{}_nA'\longrightarrow G_m
\]
et par passage à la limite
\[
T_\ell(A)\times T_\ell(A')\longrightarrow T_\ell G_m)\overset{\text{déf}}{=}\unsure{\mathbb Z|\mathbf Z}_\ell(1).
\]

\uline{\unsure{Prop.|Rmq}} Dans cet accouplement, $D(\unsure{M'_n|M_n})$ est orthogonal à ${}_nG'$, \deleted{\unsure{et}} (a un sens si ${}_nM'=0$) et \deleted{$\unsure{D(M_n)}$} ${}_nG$ est orth. à $D(M_n)$ (a un sens si ${}_nM=0$). Les accouplements induits
\figure{\begin{tikzpicture}\draw (0,0)--(0,1.8);\end{tikzpicture}}

\diagram{\begin{tikzcd}[ampersand replacement=\&, column sep=small,row sep=small]
\deleted{\unsure{{}_nA/}} \& {}_nA'/{}_nG' \& \times \& D(M'_n) \arrow[rr] \arrow[ll,bend right=75,no head] \& \& G_m\\
 \& \wr \& \& \Vert \& \&\\
 \& M'_n \& \& {}_nT \& \&
\end{tikzcd}}

et
\diagram{\begin{tikzcd}[ampersand replacement=\&, column sep=small,row sep=small]
D(M_n) \& \times \& \unsure{{}_nA|A}/{}_nG \arrow[r,no head] \& G_m\\
\Vert \& \& \wr \&\\
{}_nT \& \& M_n \&
\end{tikzcd}},

ainsi que
\diagram{\begin{tikzcd}[ampersand replacement=\&, column sep=small,row sep=small]
\deleted{\unsure{et}} \& {}_nG/{}_nT \& \times \& {}_nG'/{}_nT'\\
 \& \Vert\ \deleted{\unsure{\text{et}}} \& \& \wr\\
 \& {}_nB \& \& {}_nB'
\end{tikzcd}}
sont ceux qu'on devine.

\uline{Corollaire} Si l'accouplement entre ${}_nB$ et ${}_nB'$ est une dualité parfaite (car $\unsure{{}_nT_n}=\unsure{{}_nT'_n}=0$) \deleted{\unsure{donc}} et si ${}_nM=0$, ${}_nM'=0$, $nB=B$, $nB'=B'$, alors l'accouplement de ${}_nA$ et ${}_nA'$ est une dualité parfaite.
\figure{\begin{tikzpicture}\draw (0,0)--(0,2);\end{tikzpicture}}
\margin{32}
